% GENERATED FILE. Edit canonical lessons, then run npm run generate:books.
% canonical-source-hash: fnv1a64:548c94245f9ae1ee
% canonical-lessons: KA-C335-sampadaka, KA-C335-lekhana, KA-C335-mukhaputa, KA-C335-karapatra, KA-C335-prakasaka

\chapter{Making a Newspaper}
\label{ch:ka-making-a-newspaper}
\hlchaptermodality{335}

\begin{chapteropening}
\textbf{By the end of this chapter:} \emph{I can say an editor, an article, the front page, a leaflet and a publisher.}

Complete the last lesson of chapter 335: I can say an editor, an article, the front page, a leaflet and a publisher.
\end{chapteropening}

\section[\texorpdfstring{\kn{ಸಂಪಾದಕ} (sampādaka)}{ಸಂಪಾದಕ (sampādaka)}]{\kn{ಸಂಪಾದಕ} (sampādaka) — an editor}

\label{lesson:KA-C335-sampadaka}



\begin{quote}
Before the new one: say the Kannada for a baby shower, then the Kannada for a successor.
\end{quote}

\subsection*{You'll want to know: \kn{ಸಂಪಾದಕ}}
\textbf{\kn{ಸಂಪಾದಕ}} — \emph{sampādaka} — \textquotedblleft{}an editor\textquotedblright{}.

\begin{quote}
\kn{ಸಂಪಾದಕರು} \kn{ಆ} \kn{ಸುದ್ದಿಯನ್ನು} \kn{ಓದಿ} \kn{ಒಪ್ಪಿದರು}.
\end{quote}

\begin{quote}
\emph{sampādakaru ā suddiyannu ōdi oppidaru.}
\end{quote}

\begin{quote}
The editor read that news report and approved it.
\end{quote}

Plural: \kn{ಸಂಪಾದಕರು}, \emph{sampādakaru}.

\kn{ಸಂಪಾದಿಸು} is to edit, and also to earn.

\begin{cousinweb}
\textbf{\kn{ಸಂಪಾದಕ}} (\emph{sampādaka}) is borrowed from Sanskrit \emph{sampādaka}, \textquotedblleft{}procuring, bringing about\textquotedblright{} (Monier-Williams); the sense \textquotedblleft{}an editor\textquotedblright{}, one who puts a work together, is modern.
\end{cousinweb}

\begin{grammarlens}[title={What you've built}]
One more word for telling what happened.
\end{grammarlens}

\subsection*{Your turn}
\begin{itemize}
\raggedright
  \item \emph{Say it:} \emph{sampādaka}
  \item \emph{Say it:} \emph{sampādaka}, once more
  \item \emph{Say it:} \emph{sampādakaru ā suddiyannu ōdi oppidaru.} — The editor read that news report and approved it.
  \item \emph{Recall:} say the Kannada for a baby shower, then the Kannada for a successor, then say \emph{sampādaka} again
\end{itemize}

\begin{tcolorbox}[breakable,colback=teal!4,colframe=teal!35!black,title={Before you move on}]
What does \kn{ಸಂಪಾದಕ} mean? (\textquotedblleft{}An editor\textquotedblright{}.) Say it once more.
\end{tcolorbox}
\section[\texorpdfstring{\kn{ಲೇಖನ} (lēkhana)}{ಲೇಖನ (lēkhana)}]{\kn{ಲೇಖನ} (lēkhana) — an article; a piece of writing}

\label{lesson:KA-C335-lekhana}



\begin{quote}
Before the new one: say the Kannada for a successor, then the Kannada for an editor.
\end{quote}

\subsection*{You'll want to know: \kn{ಲೇಖನ}}
\textbf{\kn{ಲೇಖನ}} — \emph{lēkhana} — \textquotedblleft{}an article; a piece of writing\textquotedblright{}.

\begin{quote}
\kn{ಅವಳ} \kn{ಲೇಖನ} \kn{ಭಾನುವಾರದ} \kn{ಪತ್ರಿಕೆಯಲ್ಲಿ} \kn{ಬಂತು}.
\end{quote}

\begin{quote}
\emph{avaḷa lēkhana bhānuvārada patrikeyalli bantu.}
\end{quote}

\begin{quote}
Her article came out in the Sunday paper.
\end{quote}

Plural: \kn{ಲೇಖನಗಳು}, \emph{lēkhanagaḷu}.

Not to be confused with \kn{ಲೇಖನಿ}, a pen. \kn{ಲೇಖಕ} is a writer.

\begin{cousinweb}
\textbf{\kn{ಲೇಖನ}} (\emph{lēkhana}) is borrowed from Sanskrit \emph{lekhana}, \textquotedblleft{}scratching, writing\textquotedblright{} (Monier-Williams), from the root \emph{likh}, \textquotedblleft{}to scratch, to write\textquotedblright{}.
\end{cousinweb}

\begin{grammarlens}[title={What you've built}]
One more word for telling what happened.
\end{grammarlens}

\subsection*{Your turn}
\begin{itemize}
\raggedright
  \item \emph{Say it:} \emph{lēkhana}
  \item \emph{Say it:} \emph{lēkhana}, once more
  \item \emph{Say it:} \emph{avaḷa lēkhana bhānuvārada patrikeyalli bantu.} — Her article came out in the Sunday paper.
  \item \emph{Recall:} say the Kannada for a successor, then the Kannada for an editor, then say \emph{lēkhana} again
\end{itemize}

\begin{tcolorbox}[breakable,colback=teal!4,colframe=teal!35!black,title={Before you move on}]
What does \kn{ಲೇಖನ} mean? (\textquotedblleft{}An article; a piece of writing\textquotedblright{}.) Say it once more.
\end{tcolorbox}
\section[\texorpdfstring{\kn{ಮುಖಪುಟ} (mukhapuṭa)}{ಮುಖಪುಟ (mukhapuṭa)}]{\kn{ಮುಖಪುಟ} (mukhapuṭa) — the front page}

\label{lesson:KA-C335-mukhaputa}



\begin{quote}
Before the new one: say the Kannada for an editor, then the Kannada for an article.
\end{quote}

\subsection*{You'll want to know: \kn{ಮುಖಪುಟ}}
\textbf{\kn{ಮುಖಪುಟ}} — \emph{mukhapuṭa} — \textquotedblleft{}the front page\textquotedblright{}.

\begin{quote}
\kn{ಆ} \kn{ಅಪಘಾತದ} \kn{ಸುದ್ದಿ} \kn{ಮುಖಪುಟದಲ್ಲಿ} \kn{ಇತ್ತು}.
\end{quote}

\begin{quote}
\emph{ā apaghātada suddi mukhapuṭadalli ittu.}
\end{quote}

\begin{quote}
The news of that accident was on the front page.
\end{quote}

\kn{ಮುಖ}, face, and \kn{ಪುಟ}, a page: the face page. A website has one too.

\begin{cousinweb}
\textbf{\kn{ಮುಖಪುಟ}} (\emph{mukhapuṭa}) is a Kannada compound built from two Sanskrit words: \emph{mukha}, \textquotedblleft{}face, front\textquotedblright{} (CDIAL 10158), and \emph{puṭa}, \textquotedblleft{}a fold\textquotedblright{} (CDIAL 8253), which Kannada uses for \textquotedblleft{}a page\textquotedblright{}.
\end{cousinweb}

\begin{grammarlens}[title={What you've built}]
One more word for telling what happened.
\end{grammarlens}

\subsection*{Your turn}
\begin{itemize}
\raggedright
  \item \emph{Say it:} \emph{mukhapuṭa}
  \item \emph{Say it:} \emph{mukhapuṭa}, once more
  \item \emph{Say it:} \emph{ā apaghātada suddi mukhapuṭadalli ittu.} — The news of that accident was on the front page.
  \item \emph{Recall:} say the Kannada for an editor, then the Kannada for an article, then say \emph{mukhapuṭa} again
\end{itemize}

\begin{tcolorbox}[breakable,colback=teal!4,colframe=teal!35!black,title={Before you move on}]
What does \kn{ಮುಖಪುಟ} mean? (\textquotedblleft{}The front page\textquotedblright{}.) Say it once more.
\end{tcolorbox}
\section[\texorpdfstring{\kn{ಕರಪತ್ರ} (karapatra)}{ಕರಪತ್ರ (karapatra)}]{\kn{ಕರಪತ್ರ} (karapatra) — a leaflet; a pamphlet}

\label{lesson:KA-C335-karapatra}



\begin{quote}
Before the new one: say the Kannada for an article, then the Kannada for the front page.
\end{quote}

\subsection*{You'll want to know: \kn{ಕರಪತ್ರ}}
\textbf{\kn{ಕರಪತ್ರ}} — \emph{karapatra} — \textquotedblleft{}a leaflet; a pamphlet\textquotedblright{}.

\begin{quote}
\kn{ಹುಡುಗರು} \kn{ಬಸ್} \kn{ನಿಲ್ದಾಣದಲ್ಲಿ} \kn{ಕರಪತ್ರ} \kn{ಹಂಚಿದರು}.
\end{quote}

\begin{quote}
\emph{huḍugaru bas nildāṇadalli karapatra hañcidaru.}
\end{quote}

\begin{quote}
The boys handed out leaflets at the bus stand.
\end{quote}

Plural: \kn{ಕರಪತ್ರಗಳು}, \emph{karapatragaḷu}.

\kn{ಕರ}, hand, and \kn{ಪತ್ರ}, a leaf: a sheet put into people's hands.

\begin{cousinweb}
\textbf{\kn{ಕರಪತ್ರ}} (\emph{karapatra}) is borrowed from Sanskrit, a modern compound of \emph{kara}, \textquotedblleft{}the hand\textquotedblright{} (Monier-Williams), and \emph{pattra}, \textquotedblleft{}a leaf, the leaf of a book\textquotedblright{} (CDIAL 7733).
\end{cousinweb}

\begin{grammarlens}[title={What you've built}]
One more word for telling what happened.
\end{grammarlens}

\subsection*{Your turn}
\begin{itemize}
\raggedright
  \item \emph{Say it:} \emph{karapatra}
  \item \emph{Say it:} \emph{karapatra}, once more
  \item \emph{Say it:} \emph{huḍugaru bas nildāṇadalli karapatra hañcidaru.} — The boys handed out leaflets at the bus stand.
  \item \emph{Recall:} say the Kannada for an article, then the Kannada for the front page, then say \emph{karapatra} again
\end{itemize}

\begin{tcolorbox}[breakable,colback=teal!4,colframe=teal!35!black,title={Before you move on}]
What does \kn{ಕರಪತ್ರ} mean? (\textquotedblleft{}A leaflet; a pamphlet\textquotedblright{}.) Say it once more.
\end{tcolorbox}
\section[\texorpdfstring{\kn{ಪ್ರಕಾಶಕ} (prakāśaka)}{ಪ್ರಕಾಶಕ (prakāśaka)}]{\kn{ಪ್ರಕಾಶಕ} (prakāśaka) — a publisher}

\label{lesson:KA-C335-prakasaka}



\begin{quote}
Before the new one: say the Kannada for the front page, then the Kannada for a leaflet.
\end{quote}

\subsection*{You'll want to know: \kn{ಪ್ರಕಾಶಕ}}
\textbf{\kn{ಪ್ರಕಾಶಕ}} — \emph{prakāśaka} — \textquotedblleft{}a publisher\textquotedblright{}.

\begin{quote}
\kn{ಪ್ರಕಾಶಕರು} \kn{ಅಜ್ಜನ} \kn{ಕಥೆಗಳನ್ನು} \kn{ಪುಸ್ತಕವಾಗಿ} \kn{ಪ್ರಕಟಿಸಿದರು}.
\end{quote}

\begin{quote}
\emph{prakāśakaru ajjana kathegaḷannu pustakavāgi prakaṭisidaru.}
\end{quote}

\begin{quote}
The publisher brought out Grandfather's stories as a book.
\end{quote}

Plural: \kn{ಪ್ರಕಾಶಕರು}, \emph{prakāśakaru}.

\kn{ಪ್ರಕಾಶ} is light: a publisher brings a work into the light.

\begin{cousinweb}
\textbf{\kn{ಪ್ರಕಾಶಕ}} (\emph{prakāśaka}) is borrowed from Sanskrit \emph{prakāśaka}, \textquotedblleft{}bright, shining, making known\textquotedblright{} (Monier-Williams), from \emph{prakāśa}, \textquotedblleft{}light, publicity\textquotedblright{}.
\end{cousinweb}

\begin{grammarlens}[title={What you've built}]
One more word for telling what happened.
\end{grammarlens}

\subsection*{Your turn}
\begin{itemize}
\raggedright
  \item \emph{Say it:} \emph{prakāśaka}
  \item \emph{Say it:} \emph{prakāśaka}, once more
  \item \emph{Say it:} \emph{prakāśakaru ajjana kathegaḷannu pustakavāgi prakaṭisidaru.} — The publisher brought out Grandfather's stories as a book.
  \item \emph{Recall:} say the Kannada for the front page, then the Kannada for a leaflet, then say \emph{prakāśaka} again
\end{itemize}

\begin{tcolorbox}[breakable,colback=teal!4,colframe=teal!35!black,title={Before you move on}]
What does \kn{ಪ್ರಕಾಶಕ} mean? (\textquotedblleft{}A publisher\textquotedblright{}.) Say it once more.
\end{tcolorbox}
